A realização de um sistema de monitoramento para cuidado dos camarões é de grande importância, ajuda tanto o produtor a vender um produto de qualidade, tanto o consumidor que irá consumir algo de qualidade sem prejudicar a sua saúde. Por conta desta importância, é de se imaginar que existam sistemas semelhantes. Dentre os projetos já existentes, há de se destacar determinados artigos.

O trabalho de \cite{dantas2020} propõe um sistema de monitoramento da qualidade da água para a carcinicultura baseado em tecnologias de IoT, utilizando microcontroladores ESP32 integrados a módulos Longo Alcance (do inglês, Long Range, ou LoRa) e sensores de variáveis como pH, oxigênio dissolvido e temperatura. O sistema realiza coleta, processamento e transmissão dos dados via Rede de Longo Alcance e Ampla Cobertura (do inglês, Long Range Wide Area Network, ou LoRaWAN), com integração nos Serviços de Nuvem da Amazon (do inglês, Amazon Web Services, ou AWS) por meio do Protocolo de Telemetria por Enfileiramento de Mensagens (do inglês, Message Queuing Telemetry Transport, ou MQTT). Os resultados mostraram medições precisas de temperatura e pH, com variações mínimas de 0.5°C e comunicação LoRa alcançando cerca de 138 metros, o que se mostrou suficiente para pequenas e médias fazendas. Entre as limitações, destacam-se a necessidade de ampliar o alcance da rede, aprimorar a autonomia energética e incluir novos sensores para um monitoramento mais completo.

O trabalho de \cite{Uddin2020} apresenta um sistema de monitoramento inteligente para fazendas de camarão de água doce baseado em IoT, utilizando sensores conectados a um microcontrolador ESP32 para acompanhar parâmetros como temperatura, pH, salinidade, oxigênio dissolvido e turbidez. Os dados são transmitidos via Wi-Fi para a nuvem e analisados com suporte à decisão baseado na técnica de Processo de Análise Hierárquica (do inglês, Analytic Hierarchy Process, ou AHP) que auxilia na avaliação da qualidade da água e na tomada de decisões pelos produtores. Os testes mostraram alta precisão, sendo eles acima de 93\% e impacto positivo na produtividade, com aumento do tamanho, peso e taxa de sobrevivência dos camarões. Entre as limitações estão a dependência de Wi-Fi, necessidade de calibração periódica e desafios de escalabilidade.

O trabalho de \cite{Zainuddin2019} propõe um sistema de monitoramento da qualidade da água em viveiros de camarão Vannamei utilizando uma rede de sensores sem fio (do inglês, Wireless Sensor Network, ou WSN) integrada à tecnologia IoT. O sistema emprega sensores de temperatura, pH e turbidez conectados a microcontroladores Arduino Atmega 2560 e 328, além de módulos Xbee para comunicação sem fio. Os dados coletados são transmitidos por meio do módulo ESP8266 para uma aplicação web, permitindo o acompanhamento remoto das condições da água. Os testes de calibração demonstraram alta precisão, com índices acima de 97\%, comprovando a eficiência e confiabilidade do sistema. Entretanto, o desempenho depende da estabilidade da conexão de internet e apresenta limitações quanto à expansão para novos sensores e ausência de mecanismos de redundância, o que pode afetar sua escalabilidade e robustez.

Os trabalhos analisados monitoram variáveis como temperatura e pH da água, demonstrando a eficiência da IoT no acompanhamento remoto dos viveiros, mas ainda enfrentam limitações de conectividade, escalabilidade e necessidade de calibração periódica.

Em comparação, o projeto proposto apresenta uma abordagem mais completa, integrando o monitoramento de pH, temperatura e amônia ao controle automatizado da alimentação, com definição de horários e quantidades ideais de ração. Além disso, gera relatórios e gráficos históricos para análise da qualidade da água e desempenho dos cativeiros, e envia alertas preventivos em caso de variações críticas. Assim, diferencia-se por unificar monitoramento ambiental, controle alimentar e apoio à gestão do cultivo em uma única solução.